\section{Discussion}

\subsection{Interpretation}

The results presented here establish EDT as a unifying framework for understanding how quantum resources become physically usable~\cite{Horodecki2009,Streltsov2017}. Across data-driven analysis, minimal dynamical models, and sensing-oriented control optimization, a consistent picture emerges: useful quantum dynamics arise not from maximal entanglement or unrestricted Hilbert-space exploration, but from a finite regime in which structure, interaction, topology, and dynamics are jointly organized~\cite{Amico2008,Eisert2010}.

A central implication is that quantum advantage should be understood as a problem of resource organization rather than resource amplification~\cite{Preskill2018}. In conventional approaches, increasing interaction strength, circuit depth, or entanglement is often assumed to improve performance. However, the results here show that beyond a finite operating window, dynamics become overdriven, leading to degradation of coherence and correlation. This reveals a fundamental trade-off: the same interactions that generate quantum resources also destabilize them when pushed beyond the structured regime~\cite{Altman2015,Swingle2018}.

The data-driven identification of EDT in K$_3$C$_{60}$ demonstrates that this regime arises directly from microscopic electronic structure~\cite{Gunnarsson1997,Capone2009}. The existence of a finite curvature-defined window $\Delta k_{\mathrm{EDT}}$ shows that non-adiabatic coupling naturally organizes quantum resources into localized regions of momentum space~\cite{Mead1979,Yarkony1996}. This provides a physical origin for structured quantum dynamics. The supplementary archived comparison materials show that the same paired-branch residual scan can be exported to materially different datasets, while also making clear that the strength of the finite-window claim depends on branch provenance and robustness information for the specific material under study.

From a dynamical perspective, the minimal plaquette model shows that the EDT regime is characterized by the coexistence of operator spreading and persistent many-body correlations~\cite{Choi2016,KimHuse2013}. This distinguishes EDT from both weakly interacting systems, where dynamics remain local, and strongly driven systems, where coherence is rapidly lost. Instead, EDT occupies an intermediate regime in which interference is stabilized by topology and constrained by symmetry, leading to dynamically organized entanglement~\cite{Swingle2018}.

The scaling and implementation results further demonstrate that this regime is operationally robust. The persistence of the optimal loop strength $\lambda^\star$ and the EDT score under parameter variation and approximate Trotterization indicates that the regime is not fine-tuned~\cite{Lloyd1996,Childs2019}. This establishes EDT as a practical design principle for quantum technologies, guiding the selection of operating conditions in both digital and analog platforms.

The NV-center sensing results provide a sensing-oriented manifestation of the same principle~\cite{Degen2017,Barry2020}. In NV--XY8 protocols, performance is governed by competing resource factors, including filter gain, coherence time, pulse fidelity, and readout contrast. Optimal sensitivity is achieved within a finite control regime in which these factors are balanced, rather than maximized individually. This behavior mirrors the EDT regime identified in the plaquette model, demonstrating that the same structured trade-offs govern both abstract quantum dynamics and experimentally relevant control landscapes~\cite{Casola2018}.

\paragraph*{Relation to superconductivity and electron--phonon theory.}

The curvature-based origin of EDT connects directly to the microscopic theory of superconductivity. In the Bardeen--Pines formulation, electron--phonon interactions renormalize the electronic potential through dynamical screening, producing an effective attractive interaction that underlies BCS pairing~\cite{BardeenPines1955,BCS1957}. These interactions arise from the same dependence of electronic energies on lattice coordinates that generates the non-adiabatic coupling fields considered here.

From the EDT perspective, the key quantity is the second-order response of electronic energies to normal-mode coordinates, which captures curvature induced by lattice dynamics. While in conventional superconductivity this curvature contributes to pairing amplitudes, EDT interprets it more broadly as a generator of structured quantum resource manifolds. Regions of enhanced curvature then define finite domains in parameter space where electronic states are strongly coupled, interference is stabilized, and correlated dynamics are supported~\cite{Giustino2017}.

\paragraph*{Relation to Quantum Advantage Paradigms}

Early demonstrations of quantum advantage, such as random-circuit sampling on programmable superconducting processors~\cite{Arute2019}, emphasized the role of large-scale Hilbert-space exploration in generating classically intractable dynamics. In such systems, advantage arises from highly entangled, strongly scrambled states that are difficult to simulate but lack intrinsic structure or direct operational relevance.

In contrast, EDT identifies a distinct regime of quantum dynamics in which entanglement, interference, and correlation are not maximized globally, but instead organized within finite, geometry-aligned manifolds. Within this regime, operator spreading remains structured and compatible with persistent correlations, enabling both controllability and interpretability. This establishes a complementary paradigm in which useful quantum behavior arises not from the size of Hilbert space explored, but from the structure of the accessible dynamical manifold.

\paragraph*{Relation to Prior Work}

The present framework connects several strands of research across quantum physics and quantum information. The role of non-adiabatic coupling and curvature in shaping electronic structure has been extensively studied in the context of electron--phonon interactions and geometric-phase effects~\cite{Mead1979,Yarkony1996,Giustino2017}. In parallel, quantum-information approaches have explored structured representations of many-body states, including tensor-network methods and symmetry-resolved constructions such as Schur transforms and Clebsch--Gordan decompositions~\cite{Schollwock2011,Bacon2006}. EDT links these directions by identifying a physically grounded regime in which microscopic interactions generate structured manifolds that simultaneously organize entanglement, constrain dynamics, and enable efficient representations.

\subsection{Conclusion and Outlook}

We have introduced EDT as a framework in which quantum resources are organized into finite, geometry-aligned manifolds. By combining data-driven analysis of electronic structure with minimal dynamical models and implementation-level scaling, we have shown that useful quantum dynamics arise within a finite regime defined by the interplay of symmetry, interaction, topology, and dynamics.

This regime is both physically grounded and operationally robust. It emerges from microscopic non-adiabatic coupling~\cite{Mead1979,Yarkony1996}, manifests as structured operator dynamics~\cite{Swingle2018}, and remains stable under realistic circuit approximations~\cite{Lloyd1996,Childs2019}. Across materials, dynamical models, and sensing platforms, a consistent picture emerges: quantum functionality is governed by the structure of the accessible dynamical manifold rather than by the size of Hilbert space~\cite{Amico2008,Eisert2010}.

Looking forward, several directions naturally follow. On the materials side, curvature-based identification of EDT provides a route to screening structured quantum regimes in correlated and strongly coupled systems, provided that branch selection and robustness can be documented with the same care as in the K$_3$C$_{60}$ reference case~\cite{Capone2009}. On the algorithmic side, EDT motivates quantum-computation strategies based on navigating structured dynamical manifolds, rather than maximizing circuit depth or entanglement~\cite{Preskill2018}. On the sensing side, platforms such as NV centers provide a practical route for testing finite-regime control optimization under realistic constraints~\cite{Degen2017}.
